% ----------------------------------------------------------------------
%                   LATEX TEMPLATE FOR PhD THESIS
% ----------------------------------------------------------------------

% loosely based on on Harish Bhanderi's PhD/MPhil template, then Uni Cambridge
% http://www-h.eng.cam.ac.uk/help/tpl/textprocessing/ThesisStyle/
% corrected and extended in 2007 by Jakob Suckale
% perfected in 2023 by G.J. van der Gun, then MSc Human Movement Science

% Style file
% Most style definitions are in the external file PhDthesisPSnPDF.
% In this template package, it can be found in ./Latex/classes/
\documentclass[openany,11pt]{Latex/classes/PhDthesisPSnPDF}
% option 'twoside' for book style printing, i.e. start chapters only on odd pages
% option 'openany' for starting chapters on any page

% loaded packages; add or remove as required
% \usepackage{pdfpages} % simple inclusion of pdf pages
% \usepackage[numbered, framed, autolinebreaks]{mcode} % include matlab code
% \usepackage{epigraph} % nice quotes
\usepackage{soul} % highlighting for editting
\usepackage{lipsum} % generate fake latin
\usepackage{amsmath} % mathematical symbols
\usepackage{nomencl} \makenomenclature % required for glossary
\usepackage{graphicx,placeins,rotating,caption,subcaption} % include (multiple) graphics, extra float options (float barrier) and sideways figures
\usepackage{booktabs,multirow} % extra options for tables
\newcommand{\ra}[1]{\renewcommand{\arraystretch}{#1}} % make command to change default row height

\renewcommand{\floatpagefraction}{.8} % only pages with 80% floats become pure float page -> reduces amount of float-only pages and the vert. distance between floats on float-only pages.

%: ----------------------------------------------------------------------
%:                  TITLE PAGE: name, degree,..
% ----------------------------------------------------------------------

% set page layout
%\textwidth 15cm
%\textheight 22cm
%\parindent 10pt
\oddsidemargin 0.85cm \evensidemargin 0.37cm % uneven page margins make for better booklet readability

%\includeonly{sections/intro} % only compile certain chapters

\begin{document}
	
	%Macro file for Latex
	% Macros help you summarise frequently repeated Latex commands.
	% Define macros in the macroe file ./Latex/macros/MacroFile1.tex
	% An macro that you may use frequently is the figuremacro
	% This file contains macros that can be called up from connected TeX files
% It helps to summarise repeated code, e.g. figure insertion (see below).

%: ----------------------- FigureMacros ------------------------

% insert a centered figure with caption and description
% parameters 1:filename, 2:label, 3:description
\newcommand{\figuremacro}[4]{
	\begin{figure}[htbp]
		\centering
		\includegraphics[width=1\textwidth]{figures/#1}
		\caption[#3]{#4}
		\label{fig:#2}
	\end{figure}
}

% insert a centered figure with caption and description AND WIDTH
% parameters 1:filename, 2:label, 3:description, 4: textwidth
% textwidth 1 means as text, 0.5 means half the width of the text
\newcommand{\figuremacroW}[5]{
	\begin{figure}[htbp]
		\centering
		\includegraphics[width=#5\textwidth]{figures/#1}
		\caption[#3]{#4}
		\label{fig:#2}
	\end{figure}
}

% insert a centered multi figure with caption and description
% parameters 1:filename A, 2:filename B, 3:label, 4:description,
% textwidth 1 means as text, 0.5 means half the width of the text
\newcommand{\figuremacroG}[5]{
	\begin{figure*}[th]
		\centering
		\begin{subfigure}[t]{0.28\textwidth}
			\includegraphics[width=\linewidth]{figures/#1}
			\caption{}\label{fig:#1}
		\end{subfigure}
		\begin{subfigure}[t]{0.28\textwidth}
			\includegraphics[width=\linewidth]{figures/#2}
			\caption{}\label{fig:#2}
		\end{subfigure}
		\caption[#4]{#5}\label{fig:#3}
	\end{figure*}
}
	
	\frontmatter % declares frontmatter formatting
	
	\thispagestyle{empty} % removes page numbering
	
	\begin{center}
		
		Vrije Universiteit Amsterdam \hspace*{2cm} ~~Your Organisation/Company
		
		\vspace{1mm}
		
		\hspace*{-6.5cm}\includegraphics[height=20mm]{frontmatter/figures/vu-griffioen.pdf}
		
		\vspace*{-2cm}\hspace*{7.5cm}\includegraphics[height=20mm]{frontmatter/figures/your_logo.jpg}
		
		\vspace{2cm}
	
		{\Large Master Thesis}
		
		\vspace*{1.5cm}
		
		\rule{.9\linewidth}{.6pt}\\[0.4cm]
		{\huge \bfseries Title of the Thesis\par}\vspace{0.4cm}
		\rule{.9\linewidth}{.6pt}\\[1.5cm]
		
		\vspace*{2mm}
		
		{\Large
			\begin{tabular}{l}
				{\bf Author:} ~~student name ~~~~ (student number)
			\end{tabular}
		}
		
		\vspace*{2cm}
		
		\begin{tabular}{ll}
			{\it 1st supervisor:}   & ~~supervisor name \\
			{\it daily supervisor:} & ~~supervisor name ~~~~ (company, if applicable) \\
			{\it 2nd reader:}       & ~~supervisor name
		\end{tabular}
		
		\vspace*{2.5cm}
		
		\textit{A thesis submitted in fulfillment of the requirements for\\ the joint UvA-VU Master of Science degree in Computer Science}
		
		\vspace*{1.8cm}
		
		\today\\[4cm] % Date
		
	\end{center}
	
	\newpage
	
	
	% ----------------------------------------------------------------------
	
	% turn off overfull and underfull hboxes
	\hbadness=10000
	\hfuzz=50pt

	%: --------------------------------------------------------------
	%:                  FRONT MATTER: dedications, abstract,..
	% --------------------------------------------------------------
	
	%: ----------------------- generate cover page ------------------------
	
	\begin{center}
		\textit{``I am the master of my fate, I am the captain of my soul'' \\ from {\em Invictus}, by William Ernest Henley}
	\end{center}
	
	%: ----------------------- cover page back side ------------------------
	% Your research institution may require reviewer names, etc.
	
	\newpage
	%\vspace{10mm}
	%1. First Reader: Name Surname
	%
	%\vspace{10mm}
	%2. Daily Supervisor: Name Surname
	%
	%\vspace{10mm}
	%3. Second Reader: Name Surname
	%
	%\vspace{10mm}
	%4. Industrial Supervisor: Name Surname
	%
	%\vspace{20mm}
	%Day of the defense:
	
	%\vspace{20mm}
	%\hspace{70mm}Signature from head of PhD committee:
	
	
	%: ----------------------- abstract ------------------------
	
	% Your institution may have specific regulations if you need an abstract and where it is to be placed in the document. The default here is just after title.
	
	
% Thesis Abstract -----------------------------------------------------


%\begin{abstractslong}    %uncommenting this line, gives a different abstract heading
\begin{abstracts}        %this creates the heading for the abstract page
	\thispagestyle{plain} % uncomment if page numbering is required
	
	\noindent \textit{Objective}. 
	
	\noindent \textit{Methods}. 

	\noindent \textit{Results}. 
	
	\noindent \textit{Conclusions}. 
	
\end{abstracts}
%\end{abstractlongs}


% ---------------------------------------------------------------------- 

	
	%: ----------------------- contents ------------------------
	
	\setcounter{secnumdepth}{3} % section numbering depth
	\setcounter{tocdepth}{3}    % sections levels shown in ToC
	
	\tableofcontents            % print the table of contents
	% levels are: 0 - chapter, 1 - section, 2 - subsection, 3 - subsection
	
	%: ----------------------- list of figures/tables ------------------------
	
	\listoffigures	% print list of figures
	
	\listoftables % print list of tables
	
	%: ----------------------- glossary ------------------------
	
	% Tie in external source file for definitions: /frontmatter/glossary.tex
	% Glossary entries can also be defined in the main text. See glossary.tex
	% 
	%	% this file is called up by thesis.tex
% content in this file will be fed into the main document

% Glossary entries are defined with the command \nomenclature{1}{2}
% 1 = Entry name, e.g. abbreviation; 2 = Explanation
% You can place all explanations in this separate file or declare them in the middle of the text. Either way they will be collected in the glossary.

% required to print nomenclature name to page header
\markboth{\MakeUppercase{\nomname}}{\MakeUppercase{\nomname}}


% ----------------------- contents from here ------------------------%


\nomenclature{LSY}{ehbfuefebbfbjkjkebfjbfbfw} 
\nomenclature{DEPC}{diethyl-pyro-carbonate; used to remove RNA-degrading enzymes (RNAases) from water and laboratory utensils}
\nomenclature{DMSO}{dimethyl sulfoxide; organic solvent, readily passes through skin, cryoprotectant in cell culture}
\nomenclature{EDTA}{Ethylene-diamine-tetraacetic acid; a chelating (two-pronged) molecule used to sequester most divalent (or trivalent) metal ions, such as calcium (Ca$^{2+}$) and magnesium (Mg$^{2+}$), copper (Cu$^{2+}$), or iron (Fe$^{2+}$ / Fe$^{3+}$)} % all glossary entries are stored here
	%	
	%	\begin{multicols}{2} % \begin{multicols}{#columns}[header text][space]
	%	\begin{footnotesize} % scriptsize(7) < footnotesize(8) < small (9) < normal (10)
	%	
	%	\printnomenclature[1.5cm] % [] = distance between entry and description
	%	\label{nom} % target name for links to glossary
	%	
	%	\end{footnotesize}
	%	\end{multicols}
	
	%: --------------------------------------------------------------
	%:                  MAIN DOCUMENT SECTION
	% --------------------------------------------------------------
	
	% the main text starts here with the introduction, 1st chapter,...
	\mainmatter
	
	\renewcommand{\chaptername}{} % uncomment to print only "1" not "Chapter 1"
	
	
	%: ----------------------- subdocuments ------------------------
	
	% Parts of the thesis are included below. Rename the files as required.
	% But take care that the paths match. You can also change the order of appearance by moving the include commands.
	
	\chapter{Introduction}\label{s:intro} \FloatBarrier

Human movement is a complex and multi-faceted phenomenon that has been studied for decades. It encompasses a wide range of activities, from simple everyday movements such as walking and reaching, to more complex movements such as athletic performance and dance \cite{pereira2017effect,tyler2016new}. The study of human movement is an interdisciplinary field that draws on a range of scientific disciplines, including biomechanics, physiology, psychology, and neuroscience \cite{boyd2017effects}. Understanding how the body moves, how it is controlled and how it interacts with the environment is important for a wide range of applications, including sport and exercise performance, rehabilitation, and the prevention of injury\cite{van2016neuroplasticity}.

% NOTE: for explanation how to use please check MacroFile1.text
\figuremacro{example1}{ExampleLabel}{Example Title}{Example figure description}

One of the key challenges in studying human movement is the variability and complexity of the movements themselves\cite{miller2017predictors,muller2018physical,pereira2017effect}. There is a great deal of individual variation in the way people move, which can make it difficult to generalize findings across different populations. In addition, human movement is influenced by a range of factors, including physiological, psychological and environmental factors \cite{decre2016reliability,johnson2018influence}. As a result, understanding the underlying mechanisms of human movement requires a multi-disciplinary approach that incorporates a range of different research methodologies.

The aim of this thesis is to explore the factors that influence human movement and to investigate the underlying mechanisms that drive it. Specifically, the thesis will focus on the study of human movement in relation to physical activity and exercise performance. Physical activity is a key determinant of health, and understanding the mechanisms by which it affects the body can have important implications for the prevention and treatment of a range of health conditions \cite{lee2017exploring,maclellan2017quantifying,boyd2017effects}.

\figuremacroG{multi1}{multi2}{GridExampleLabel}{Grid Example Title}{Example grid figure description, with reference to Figure \ref{fig:ExampleLabel}}

The thesis will be structured into four chapters, each of which will focus on a different aspect of human movement. Chapter one will provide an overview of the key concepts and theoretical frameworks that underpin the study of human movement, including the role of biomechanics, physiology, and psychology. Chapter two will focus on the measurement of physical activity and exercise performance, including the use of wearable technology and other objective measures. Chapter three will explore the physiological and psychological mechanisms that underpin the effects of physical activity on the body, including the role of exercise intensity, duration, frequency and a random mention of \nomenclature{DMSO}. Finally, chapter four will investigate the role of environmental factors in influencing human movement, including the impact of temperature, humidity, and altitude.

% nice table style with optional footnote. Avoid vertical boundaries and bold face text. 
\begin{table}[htp]
	\centering \ra{1.2}
	\caption[Example Title]{Example table description}
	\begin{tabular}{@{}lll@{}} 
		\toprule
									& Fancy Symbols				& Some Numbers \\ 
		\midrule
		Row							& $m$						& 42\\
		Another row					& $\oint^a$ 				& 360 \\
		Row 4						& ${\zeta_z}$				& $e$ \\
		\bottomrule
		\multicolumn{3}{l}{\footnotesize{$^a$Some symbols require the Amsmath package}}
	\end{tabular}
	\label{tab:example}
\end{table}

In order to achieve these aims, a range of different research methodologies will be employed. These will include laboratory-based studies, field studies, and epidemiological studies. Laboratory-based studies will be used to investigate the physiological and psychological mechanisms that underpin the effects of physical activity on the body. Field studies will be used to explore the impact of environmental factors on human movement, while epidemiological studies will be used to investigate the relationship between physical activity and health outcomes.

Overall, this thesis aims to contribute to the growing body of knowledge on human movement and its underlying mechanisms. By exploring the factors that influence human movement in relation to physical activity and exercise performance, the thesis will provide new insights into how the body responds to different types of physical activity, and how this can be optimized for health and performance. The findings of this thesis may have important implications for the development of interventions aimed at promoting physical activity, improving exercise performance, and preventing and treating a range of health conditions.


	\chapter{Methods}\label{s:methods}\FloatBarrier

\lipsum[1-8]


	\chapter{Results}\label{s:results} \FloatBarrier

\lipsum[1-8]


	\chapter{Discussion}\label{s:discussion} \FloatBarrier

\lipsum[1-8]


	\chapter{Conclusion}\label{s:conclusion} \FloatBarrier

\lipsum[1-2]
	
	% --------------------------------------------------------------
	%:                  BACK MATTER: appendices, refs,..
	% --------------------------------------------------------------
	% the back matter: appendix and references close the thesis
	
	%: ----------------------- bibliography ------------------------
	
	\bibliographystyle{plain}
	\bibliography{./thesisbib} % adjust this to fit your BibTex file
	\addcontentsline{toc}{chapter}{References} % writes references to ToC (you may change the name here)
	
	%: ----------------------- appendix ------------------------
	
	% in alphabetical order
	\appendix
	\chapter*{Appendix} \FloatBarrier 

\lipsum[1]
	
\end{document}
